\documentclass{beamer}
\usepackage{amsmath}
\usepackage{amssymb}
\usepackage{xcolor}
\usepackage{tikz}
\usetikzlibrary{decorations.pathreplacing}

% ================================================================
% Quick Questions deck: percentages.
%
% Every question gets two slides: the question on its own for the
% mini whiteboards, then the same question again with the solution
% underneath in red. Within each topic the ten questions get
% gradually harder.
% ================================================================

\setbeamertemplate{navigation symbols}{}
\setbeamertemplate{footline}{}

% Topic divider.
\newcommand{\qtopic}[1]{%
  \section{#1}%
  \begin{frame}[plain]
    \vfill
    \begin{center}
      {\usebeamercolor[fg]{structure}\Huge\bfseries #1}
    \end{center}
    \vfill
  \end{frame}%
}

% \qq[<question size>][<solution size>]{<question>}{<solution>}
% The two optional sizes drop the type for the wordier questions and
% the longer solutions, so nothing runs off the slide.
\NewDocumentCommand{\qq}{O{\Huge}O{\LARGE}+m+m}{%
  \begin{frame}
    \vfill
    \begin{center}
      \begin{minipage}{0.92\textwidth}
        \centering #1 #3
      \end{minipage}
    \end{center}
    \vfill
  \end{frame}%
  \begin{frame}
    \vfill
    \begin{center}
      \begin{minipage}{0.92\textwidth}
        \centering #1 #3
        \par\vspace{0.9em}
        {\color{red}#2 #4}
      \end{minipage}
    \end{center}
    \vfill
  \end{frame}%
}

% Stepped working, written in the form
%   80% is 64 / 10% is 8 / 100% is 80
% Rows are separated with \ and the two columns with &.
\newcommand{\steps}[1]{%
  \begin{tabular}{@{}r@{\ is\ }l@{}}#1\end{tabular}%
}

% \barmodel{<cells>}{<cells spanned by the lower brace>}{<lower label>}{<upper label>}
% The whole bar is 9cm wide and represents the upper label; the lower
% brace picks out the part of it named by the lower label.
\newcommand{\barmodel}[4]{%
  \begin{tikzpicture}[x=1cm,y=1cm]
    \pgfmathsetmacro{\cw}{9/#1}
    \fill[red!12] (0,0) rectangle ({#2*\cw},0.7);
    \foreach \i in {1,...,#1}
      {\draw[red!60!black] ({(\i-1)*\cw},0) rectangle ({\i*\cw},0.7);}
    \draw[red!60!black,thick] (0,0) rectangle (9,0.7);
    \draw[red,decorate,decoration={brace,amplitude=4pt}] (0,0.9) -- (9,0.9)
      node[midway,above=4pt,red,font=\normalsize]{#4};
    \draw[red,decorate,decoration={brace,mirror,amplitude=4pt}] (0,-0.2) -- ({#2*\cw},-0.2)
      node[midway,below=4pt,red,font=\normalsize]{#3};
  \end{tikzpicture}%
}

% Contents-slide bullets, colour-coded by difficulty band:
% green for the first two topics, orange for the next two,
% red for the last three.
\setbeamertemplate{section in toc}{%
  \ifnum\inserttocsectionnumber<3
    \textcolor{green!55!black}{$\bullet$}%
  \else\ifnum\inserttocsectionnumber<5
    \textcolor{orange}{$\bullet$}%
  \else
    \textcolor{red}{$\bullet$}%
  \fi\fi
  ~\inserttocsection%
}

\title{Percentages: Quick Questions}
\subtitle{Mini whiteboard questions, with solutions}
\author{}
\date{}

\begin{document}

\frame{\titlepage}

\begin{frame}{Contents}
  \small
  \tableofcontents
\end{frame}

%================================================================
\qtopic{Converting between decimals and percentages}
%================================================================

\qq[\LARGE]{Write $0.5$ as a percentage}{$0.5\times 100=50\%$}
\qq[\LARGE]{Write $0.25$ as a percentage}{$0.25\times 100=25\%$}
\qq[\LARGE]{Write $30\%$ as a decimal}{$30\div 100=0.3$}
\qq[\LARGE]{Write $0.07$ as a percentage}{$0.07\times 100=7\%$}
\qq[\LARGE]{Write $6\%$ as a decimal}{$6\div 100=0.06$}
\qq[\LARGE]{Write $0.125$ as a percentage}{$0.125\times 100=12.5\%$}
\qq[\LARGE]{Write $4.5\%$ as a decimal}{$4.5\div 100=0.045$}
\qq[\LARGE]{Write $1.3$ as a percentage}{$1.3\times 100=130\%$}
\qq[\LARGE]{Write $250\%$ as a decimal}{$250\div 100=2.5$}
\qq[\LARGE]{Write $0.0025$ as a percentage}{$0.0025\times 100=0.25\%$}

%================================================================
\qtopic{Writing one number as a percentage of another}
%================================================================

\qq[\LARGE][\Large]{Write $25$ as a percentage of $100$}
  {$\dfrac{25}{100}=25\%$}
\qq[\LARGE][\Large]{Write $7$ as a percentage of $10$}
  {$\dfrac{7}{10}=\dfrac{70}{100}=70\%$}
\qq[\LARGE][\Large]{Write $13$ as a percentage of $20$}
  {$\dfrac{13}{20}=\dfrac{65}{100}=65\%$}
\qq[\LARGE][\Large]{Write $12$ as a percentage of $25$}
  {$\dfrac{12}{25}=\dfrac{48}{100}=48\%$}
\qq[\LARGE][\Large]{Write $18$ as a percentage of $40$}
  {$\dfrac{18}{40}=\dfrac{45}{100}=45\%$}
\qq[\LARGE][\Large]{Write $21$ as a percentage of $50$}
  {$\dfrac{21}{50}=\dfrac{42}{100}=42\%$}
\qq[\LARGE][\Large]{Write $9$ as a percentage of $12$}
  {$\dfrac{9}{12}=\dfrac{3}{4}=\dfrac{75}{100}=75\%$}
\qq[\LARGE][\Large]{Write $27$ as a percentage of $60$}
  {$\dfrac{27}{60}=\dfrac{9}{20}=\dfrac{45}{100}=45\%$}
\qq[\LARGE][\Large]{Write $34$ as a percentage of $80$}
  {$\dfrac{34}{80}=\dfrac{17}{40}=\dfrac{42.5}{100}=42.5\%$}
\qq[\LARGE][\Large]{In a test Ali scored $57$ out of $75$.\\[0.4em]
What percentage did Ali score?}
  {$\dfrac{57}{75}=\dfrac{19}{25}=\dfrac{76}{100}=76\%$}

%================================================================
\qtopic{Basic percentages without a calculator}
%================================================================

\qq{$50\%$ of $80$}{$80\div 2=40$}
\qq{$10\%$ of $70$}{$70\div 10=7$}
\qq{$10\%$ of $45$}{$45\div 10=4.5$}
\qq{$1\%$ of $300$}{$300\div 100=3$}
\qq{$20\%$ of $45$}{\steps{$10\%$ & $4.5$ \\ $20\%$ & $9$}}
\qq{$5\%$ of $80$}{\steps{$10\%$ & $8$ \\ $5\%$ & $4$}}
\qq{$30\%$ of $250$}{\steps{$10\%$ & $25$ \\ $30\%$ & $75$}}
\qq{$1\%$ of $42$}{$42\div 100=0.42$}
\qq{$15\%$ of $60$}{\steps{$10\%$ & $6$ \\ $5\%$ & $3$ \\ $15\%$ & $9$}}
\qq{$35\%$ of $240$}{\steps{$10\%$ & $24$ \\ $30\%$ & $72$ \\ $5\%$ & $12$ \\ $35\%$ & $84$}}

%================================================================
\qtopic{Increasing and decreasing without a calculator}
%================================================================

\qq[\LARGE]{Increase $80$ by $10\%$}
  {\steps{$10\%$ & $8$}\par\vspace{0.4em}$80+8=88$}
\qq[\LARGE]{Decrease $60$ by $50\%$}
  {\steps{$50\%$ & $30$}\par\vspace{0.4em}$60-30=30$}
\qq[\LARGE]{Increase $40$ by $25\%$}
  {\steps{$25\%$ & $10$}\par\vspace{0.4em}$40+10=50$}
\qq[\LARGE]{Decrease $90$ by $10\%$}
  {\steps{$10\%$ & $9$}\par\vspace{0.4em}$90-9=81$}
\qq[\LARGE]{Increase \pounds $250$ by $20\%$}
  {\steps{$10\%$ & \pounds $25$ \\ $20\%$ & \pounds $50$}\par\vspace{0.4em}
   \pounds $250+$ \pounds $50=$ \pounds $300$}
\qq[\LARGE]{Decrease $45$ by $20\%$}
  {\steps{$10\%$ & $4.5$ \\ $20\%$ & $9$}\par\vspace{0.4em}$45-9=36$}
\qq[\LARGE]{Increase $120$ by $5\%$}
  {\steps{$10\%$ & $12$ \\ $5\%$ & $6$}\par\vspace{0.4em}$120+6=126$}
\qq[\LARGE]{Decrease \pounds $64$ by $25\%$}
  {\steps{$25\%$ & \pounds $16$}\par\vspace{0.4em}
   \pounds $64-$ \pounds $16=$ \pounds $48$}
\qq[\LARGE]{Increase $350$ by $30\%$}
  {\steps{$10\%$ & $35$ \\ $30\%$ & $105$}\par\vspace{0.4em}$350+105=455$}
\qq[\LARGE]{Decrease $480$ by $35\%$}
  {\steps{$10\%$ & $48$ \\ $30\%$ & $144$ \\ $5\%$ & $24$ \\ $35\%$ & $168$}%
   \par\vspace{0.4em}$480-168=312$}

%================================================================
\qtopic{Finding the original value without a calculator}
%================================================================

\qq[\Large][\large]{$10\%$ of a number is $7$.\\[0.3em]What is the number?}
  {\barmodel{10}{1}{$10\%$ is $7$}{$100\%$ is $?$}\par\vspace{0.9em}
   \steps{$10\%$ & $7$ \\ $100\%$ & $70$}}

\qq[\Large][\large]{$50\%$ of a number is $24$.\\[0.3em]What is the number?}
  {\barmodel{2}{1}{$50\%$ is $24$}{$100\%$ is $?$}\par\vspace{0.9em}
   \steps{$50\%$ & $24$ \\ $100\%$ & $48$}}

\qq[\Large][\large]{$20\%$ of a number is $12$.\\[0.3em]What is the number?}
  {\barmodel{10}{2}{$20\%$ is $12$}{$100\%$ is $?$}\par\vspace{0.9em}
   \steps{$20\%$ & $12$ \\ $10\%$ & $6$ \\ $100\%$ & $60$}}

\qq[\Large][\large]{$80\%$ of a number is $64$.\\[0.3em]What is the number?}
  {\barmodel{10}{8}{$80\%$ is $64$}{$100\%$ is $?$}\par\vspace{0.9em}
   \steps{$80\%$ & $64$ \\ $10\%$ & $8$ \\ $100\%$ & $80$}}

\qq[\Large][\large]{$30\%$ of a number is $18$.\\[0.3em]What is the number?}
  {\barmodel{10}{3}{$30\%$ is $18$}{$100\%$ is $?$}\par\vspace{0.9em}
   \steps{$30\%$ & $18$ \\ $10\%$ & $6$ \\ $100\%$ & $60$}}

\qq[\Large][\large]{$25\%$ of a number is $15$.\\[0.3em]What is the number?}
  {\barmodel{4}{1}{$25\%$ is $15$}{$100\%$ is $?$}\par\vspace{0.9em}
   \steps{$25\%$ & $15$ \\ $100\%$ & $60$}}

\qq[\Large][\large]{$15\%$ of a number is $12$.\\[0.3em]What is the number?}
  {\barmodel{20}{3}{$15\%$ is $12$}{$100\%$ is $?$}\par\vspace{0.9em}
   \steps{$15\%$ & $12$ \\ $5\%$ & $4$ \\ $100\%$ & $80$}}

\qq[\Large][\large]{A coat is reduced by $20\%$ in a sale.\\[0.3em]
It now costs \pounds $48$. What was the original price?}
  {\barmodel{10}{8}{$80\%$ is \pounds $48$}{$100\%$ is $?$}\par\vspace{0.9em}
   \steps{$80\%$ & \pounds $48$ \\ $10\%$ & \pounds $6$ \\ $100\%$ & \pounds $60$}}

\qq[\Large][\large]{A price rises by $10\%$ to \pounds $88$.\\[0.3em]
What was the price before the rise?}
  {\barmodel{11}{10}{$100\%$ is $?$}{$110\%$ is \pounds $88$}\par\vspace{0.9em}
   \steps{$110\%$ & \pounds $88$ \\ $10\%$ & \pounds $8$ \\ $100\%$ & \pounds $80$}}

\qq[\Large][\large]{A house rises in value by $30\%$ to \pounds $169\,000$.\\[0.3em]
What was it worth before the rise?}
  {\barmodel{13}{10}{$100\%$ is $?$}{$130\%$ is \pounds $169\,000$}\par\vspace{0.9em}
   \steps{$130\%$ & \pounds $169\,000$ \\ $10\%$ & \pounds $13\,000$ \\ $100\%$ & \pounds $130\,000$}}

%================================================================
\qtopic{Percentage word problems}
%================================================================

\qq[\Large][\Large]{A coat costs \pounds $40$.\\[0.3em]
In a sale it is reduced by $10\%$.\\[0.3em]What is the sale price?}
  {\steps{$10\%$ & \pounds $4$}\par\vspace{0.4em}
   \pounds $40-$ \pounds $4=$ \pounds $36$}

\qq[\Large][\Large]{$30\%$ of a class of $30$ students walk to school.\\[0.3em]
How many students walk to school?}
  {\steps{$10\%$ & $3$ \\ $30\%$ & $9$}\par\vspace{0.4em}
   $9$ students}

\qq[\Large][\Large]{A phone costs \pounds $250$ plus $20\%$ VAT.\\[0.3em]
What is the total price?}
  {\steps{$10\%$ & \pounds $25$ \\ $20\%$ & \pounds $50$}\par\vspace{0.4em}
   \pounds $250+$ \pounds $50=$ \pounds $300$}

\qq[\Large][\Large]{Sam earns \pounds $800$ a month and is given a $5\%$ pay rise.\\[0.3em]
What is Sam's new monthly pay?}
  {\steps{$10\%$ & \pounds $80$ \\ $5\%$ & \pounds $40$}\par\vspace{0.4em}
   \pounds $800+$ \pounds $40=$ \pounds $840$}

\qq[\Large][\Large]{A jacket cost \pounds $60$ and now costs \pounds $45$.\\[0.3em]
What is the percentage reduction?}
  {Reduction $=$ \pounds $15$\par\vspace{0.4em}
   $\dfrac{15}{60}=\dfrac{25}{100}=25\%$}

\qq[\Large][\Large]{A town's population grows from $8000$ to $8600$.\\[0.3em]
What is the percentage increase?}
  {Increase $=600$\par\vspace{0.4em}
   $\dfrac{600}{8000}=\dfrac{7.5}{100}=7.5\%$}

\qq[\Large][\large]{A laptop costs \pounds $480$ in a sale, after a $20\%$ discount.\\[0.3em]
What was the price before the sale?}
  {\barmodel{10}{8}{$80\%$ is \pounds $480$}{$100\%$ is $?$}\par\vspace{0.9em}
   \steps{$80\%$ & \pounds $480$ \\ $10\%$ & \pounds $60$ \\ $100\%$ & \pounds $600$}}

\qq[\Large][\Large]{Maya scored $34$ out of $40$ in a test.\\[0.3em]
Ben scored $82\%$. Who did better?}
  {$\dfrac{34}{40}=\dfrac{85}{100}=85\%$\par\vspace{0.4em}
   $85\%>82\%$, so Maya did better}

\qq[\Large][\Large]{A television costs \pounds $1200$.\\[0.3em]
It is reduced by $15\%$, then by a further $10\%$ of the new price.\\[0.3em]
What is the final price?}
  {\steps{$15\%$ of \pounds $1200$ & \pounds $180$ \\
          new price & \pounds $1020$ \\
          $10\%$ of \pounds $1020$ & \pounds $102$}\par\vspace{0.4em}
   Final price $=$ \pounds $918$}

\qq[\Large][\Large]{A car costs \pounds $12\,000$.\\[0.3em]
It loses $20\%$ of its value in the first year and $15\%$ of its value in the second year.\\[0.3em]
What is it worth after two years?}
  {\steps{after year $1$ & \pounds $12\,000\times 0.8=$ \pounds $9600$ \\
          after year $2$ & \pounds $9600\times 0.85=$ \pounds $8160$}}

%================================================================
\qtopic{Harder percentages of a number}
%================================================================

\qq{$11\%$ of $200$}
  {\steps{$10\%$ & $20$ \\ $1\%$ & $2$ \\ $11\%$ & $22$}}
\qq{$21\%$ of $400$}
  {\steps{$10\%$ & $40$ \\ $20\%$ & $80$ \\ $1\%$ & $4$ \\ $21\%$ & $84$}}
\qq{$17.5\%$ of $40$}
  {\steps{$10\%$ & $4$ \\ $5\%$ & $2$ \\ $2.5\%$ & $1$ \\ $17.5\%$ & $7$}}
\qq{$12.5\%$ of $80$}
  {\steps{$10\%$ & $8$ \\ $2.5\%$ & $2$ \\ $12.5\%$ & $10$}}
\qq{$17.5\%$ of $240$}
  {\steps{$10\%$ & $24$ \\ $5\%$ & $12$ \\ $2.5\%$ & $6$ \\ $17.5\%$ & $42$}}
\qq{$32\%$ of $150$}
  {\steps{$10\%$ & $15$ \\ $30\%$ & $45$ \\ $1\%$ & $1.5$ \\ $2\%$ & $3$ \\ $32\%$ & $48$}}
\qq{$17.5\%$ of $68$}
  {\steps{$10\%$ & $6.8$ \\ $5\%$ & $3.4$ \\ $2.5\%$ & $1.7$ \\ $17.5\%$ & $11.9$}}
\qq{$87.5\%$ of $96$}
  {\steps{$100\%$ & $96$ \\ $10\%$ & $9.6$ \\ $2.5\%$ & $2.4$ \\ $12.5\%$ & $12$ \\ $87.5\%$ & $96-12=84$}}
\qq{$6.5\%$ of $400$}
  {\steps{$1\%$ & $4$ \\ $6\%$ & $24$ \\ $0.5\%$ & $2$ \\ $6.5\%$ & $26$}}
\qq[\Huge][\LARGE]{$17.5\%$ of \pounds $312$}
  {\steps{$10\%$ & \pounds $31.20$ \\ $5\%$ & \pounds $15.60$ \\
          $2.5\%$ & \pounds $7.80$ \\ $17.5\%$ & \pounds $54.60$}}

\end{document}
